\section{怎样给文言文断句和加标点}

给文言文断句和加上恰当的标点符号是检查文言文阅读能力的一个有效手段。
我们知道，
古人写的文章是没有标点符号的，
这就给初学古文的人设置了一道障碍。要跨越这道障碍，
不光要有使用标点符号的一般知识，
还要有一系列的古代汉语知识。
那么，
除了广泛学习古代汉语知识和多读文言文之外，
就没有什么方法可供参考吗？
那倒也不是。
我们只要懂得文言常用词和基本句式的一般常识，
就可以在反复通读全文，
大致了解文意的前提下，
借助一些“指路碑”，
在茫茫绊绊的文言文中找出可通的路来，
再由段到层，
由层到句，
由句到逗地缩小断句的范围，
推敲难点，
融会贯通句子的结构和句意，
分别断句和加上恰当的标点符号。
这样，
一篇文言文的各个文句就能够清晰地显示出来，
并且也说明你是基本了解这篇文言文的意思了。

这些“指路碑”指何而言呢？
大致说来，
有以下一些：
\newpage

\textbf{1. 句末助词。}

例如，
“矣、耳、焉、也”等经常用于陈述句末尾，
在这些助词之后，
一般可以考虑用句号；
“乎、耶、欤、邪”等经常用于疑问句末尾，
在这些助词之后，
一般可以考虑用问号；
“哉”字经常用于感叹句或疑问句末尾，
在它之后可以考虑用叹号或问号。
我们根据这些句末助词在句中的位置，
就能将一篇文言文点断若干处，
并能加上恰当的标点符号。

\textbf{2. 句首助词。}

例如，
“盖、夫、且、且夫、至若、若夫”
等经常用于一句话的开头，
或解释原因，
或引起议论，
或表示进层，
或表示转折，
说明情况。
我们根据这些句首助词在句中的位置，
可以考虑将这些词前面的文字用句号或逗号断开，
也可以借助它们理解下面的文意。

\textbf{3. 表示对话关系的动词。}

例如，
“曰、云”
等词之后，
一般可以点断，
加上冒号和引号。
（当然，
如果一句话当中有引用他人的话，
即使句子中有“曰、云”等词，
也不一定点断。）

\textbf{4. 介词和连词。}

例如，
“于、以”
等介词，
只能用在句子之中，
组成介词结构，
去修饰或限制动词，
绝不能在它们之后加任何标点符号；
“而、则”
等连词只能起连接作用，
在它们之后也不能加任何标点符号。
但是它们却可以指示句中词语与词语的关系，
帮助我们断句和理解文意。

\textbf{5. 表示疑问语气的词或词组。}

例如，
“何、安、胡、曷、奚、盍、焉、孰、孰与、孰若、何如、何以、何为、奈何、若何、如之何、奈之何、若之何”
等，
在这些词或词组之后，
应考虑到有构成疑问句的可能，
只要贯串上下文意，
细心推敲，
就可在句末加上问号。

\textbf{6. 表示特定结构的文言句式。}

例如：
“······者，······也。”
是典型的判断句式；
另外还有
“······，······也”，
“······者，······”，
“······者也”
等句式。
遇到这类句式，
标点符号很容易确定
（有些表示判断关系的词，
如“为、乃、即、则、皆、非”等，
也可以帮助我们断句和加标点）。

再如：
“不亦······乎”、
“得无······乎”、
“无乃······乎”、
“何以······为”、
“何······为”、
“安······哉”、
“孰与······乎”、
“与其······孰若······”
等，
是典型的反问句型。
遇到这类句式，
标点符号也极易确定。

还有
“为······所······”、
“见······于······”、
“受······于······”等，
是典型的被动句式，
可以帮助我们理解句意，
进行断句和加标点。

至于主谓倒置、
定语后置、
宾语前置等特殊句式，
我们只要掌握了它们的格式，
都可以准确地断句和加标点。

上面说的只是一般情况，
我们必须在反复体会文意，
弄清句与句的关系和句子本身的结构的基础上灵活运用。
例如“夫”字，
除了作句首助词之外，
还可用作句末助词，
放在句尾，
表示感叹语气，
如
“信夫！”、
“悲夫！”
那就得在它之后加感叹号；
它还常常用在句子之中作远指代词用，
如“独不知大狼乎？”
在这种情况下，
它的前后则不能加任何标点符号。

除了上述的一些
“指路碑”
外，
我们还要注意以下几个问题。

1. 分清
“者”、
“也”
在不同情况下的不同断句法。

由于这两个文言虚词特别频繁地出现，
而且除了判断句式而外，
情况比较复杂，
我们在断句时应对它们加倍注意。
这两个字有一条是共同的，
即只能用在文言句的句中或句末，
而不能用在句首，
因此，
在这两个字的前面，
绝不能加任何标点符号。

“者”
字用作定语后置的煞尾字时，
它的后面要点断，
用逗号。
如
“求人可使报秦者，未得。”

“者”
字用在
“有······者”
的格式中，
它的后面要点断，
用逗号。
如
“有敢为魏王使通（通报）者，死。”、
“郑人有欲买履者，先自度其足······”。

“者”
字附在别的词语之后构成
“者字结构”，
含有
“的（人、事）”
的意思，
或者用在句中表示提示，
协调语气，
这些都是句中的一个成分，
在断句处就加标点，
不在断句处就不加标点，
为了提示则必加标点。
如：
“何者？严大国之威以修敬也。”
（表提示）；
“肉食者鄙，未能远谋。”
（“者”字结构，不在断句处）；
“夺项王天下者，必沛公也。”
（“者”字结构，在断句处）；
“今者臣来，过易水，蚌方出曝。”（协调语气，不加标点）。

“也”字除用在判断句末尾之外，
还可以用在表解释、
表感叹的句子末尾，
要点断，
用句号。
如
“敌人远我，欲以火器困我也。”、
“辘辘远听，杳不知其所之也。”

“也”字用在句中，
表停顿或舒缓语气，
则视句子长短情况在
“也”
字之后用逗号点断或不加标点。
如
“顺风而呼，声非加疾也，而闻者彰。”
（表停顿）；
“师道之不传也久矣！欲人之无惑也难矣！”
（表舒缓语气）。

注意：
由于“也”字不能用于句首，
所以决不能在“也”字前加任何标点符号，
这一点要与现代汉语副词“也”字区别开来（当现代汉语“也”讲的文言副词是“亦”字）。
例如，
有位初学古文的青年同志将《淮南子·主术训》中这段文字点错了一些地方，
就是由于没有明确这一点所致。
他是这样用的标点符号，
请你帮他改正过来：

耕之为事，
也劳；
织之为事，
也扰。
扰劳之事而不舍，
者知其可以衣食也。
人之情不能无衣食，
衣食之道必始于耕织，
万民之所公见也。
物之若耕织，
者始初甚劳，
终必利，
也众。

[按：
应将首句两个“也”字之前逗号和末句“也”字之前的逗号去掉，
使“也”表示舒缓语气，
第二句和末句的“者”字之间的逗号，
都应去掉，
而在两个“者”字后面加逗号，
表示是处在句中停顿位置上的“者”字结构。]

2. 分清对话、
转述、
引用以及叙述等不同情况。
如前所述，
文中若有
“曰、云”
等表示对话的动词，
标点符号容易确定。
但是，
如果对话中又转述对话，
或者对话中承前或蒙后省略了
“曰”、
“云”
等词，
就容易混淆了。
我们应该根据上下文，
仔细体会文意，
分清哪儿是直接的对话，
哪儿是对话中的转述对话，
哪儿是引用别人的话，
而按照现代汉语标点符号的用法分别给以妥当的处理。
例如：

（孟子）曰：“独乐乐，与人乐乐，孰乐？”
（齐王）曰：“不若与人。”
曰：“与少乐乐，与众乐乐，孰乐？”
曰：“不若与众。”
臣请为王言乐。
今王鼓乐于此，
百姓闻王钟鼓之声，
管籥之音，
举疾首蹙额而相告曰：
“‘吾王之好鼓乐，夫何使我至于此极也！父子不相见，兄弟妻子离散！’······”

这一段话一律省略了主语，
在“臣请为王言乐······”
之前，
连“曰”字也省略了；
而在
“举疾首蹙额而相告曰”
之后是转引别人的话。
如果把原文的标点符号一律去掉，
要我们去加标点符号的话，
那就必须贯通文意，
辨析以上所说的情况，
而分别加上双引号或单引号，
并且要弄清对话的起止处才行。

我们如果掌握了断句和加标点符号的“指路碑”，
并且在贯通文意的前提下注意到对于对话中的省略、
转引等情况的辨析，
就能给下面所引《孟子·梁惠王上》中的一段话加上正确的标点符号：

（齐宣王）曰：“德何如则可以王矣？”（孟子）曰：“保民而王，莫之能御也。”曰：“若寡人者，可以保民乎哉？”曰：“可。”曰：“何由知吾可也？”曰：“臣闻之胡齕曰：‘王坐于堂上，有牵牛而过堂下者，王见之，曰：“牛何之？”对曰：“将以衅钟（用血涂抹钟鼓）。”王曰：“舍之！吾不忍其觳觫（恐惧的样子），若无罪而就死地。”对曰：“然则废衅钟与？”曰：“何可废也？以羊易之。”’不识有诸？”曰：“有之。”曰：“是心足以王矣。百姓皆以王为爱（吝惜）也，臣固知王之不忍也。”）

[\ 按：从“臣闻之胡齕曰”以下到“不识有诸”之前，
都是转述胡齕的话；
而胡齕所说的话中又有对话。
为了区别转述的话和转述中所说的对话，
应将转述中所说的对话一律再用“ ”表示出来，
以免混淆不清。
在“不识有诸”之前还得加上一个后单引号（’）与前面“臣闻之胡齕曰”后面的前单引号（‘）呼应，
表示转述胡齕的话到此结束。]

3. 分清专有名词（人名、地名、国名、书名、朝代、官职、年号等）和非专有名词。
一般说来，
专有名词常用作主语或宾语，
有时用作定语。
关于人的姓要注意复姓，
如“司马、欧阳、尉迟、长孙、公孙、诸亮”等，
不可拆开点断；
如果前面已出现了人的姓名，
后面可以省略姓而只称名。
例如：
“李白少读书，未成，弃去。道逢老妪磨杵（铁杵），白问其故。曰：‘欲作针。’白感其言，遂卒业。”
我们如果点成
“道逢老妪磨杵白，问其故曰······”
那就错了。
至于其他的专门名词，
我们只要在掌握一般常识的基础上，
对句子的结构和意思细心分析，
是不难弄清的。

另外，
我们还可以借助古人经常使用一些对偶句和排比句这个特点来断句和加标点符号，
对于赋体散文尤其可以这样做；
也可以借助古人写文章常常先总说后分说或者先分说后总说的特点来断句和加标点，
等等。
如果在断句或加标点时，
发现句子意思讲不通，
或者似乎能构成个句子而不符合情理，
就得细心琢磨上下文意和句子中每个词的确切含义，
加以改正。

我们只要掌握了文言文断句的“指路碑”和应该注意的事项，
再经过必要的练习，
给文言文断句和加标点符号这个问题是有可能解决好的。
我们试举一例，
以见一般：

子墨子曰万事莫贵于义今谓人曰予子冠履而断予之手足子为之乎必不为何故则冠履不若手足之贵也又曰予子天下而杀予之身子为之乎必不为何故则天下不若身之贵也争一言以相杀是贵义于其身也
\begin{flushright}
    ---《墨子·贵义》
\end{flushright}

我们在大致了解这段文意的基础上，
先找那些能帮助我们断句和加标点符号的指路碑，
可以发现：
其中有三个表示对话的动词“曰”，
则可以考虑是否该在它们之后各加冒号和引号；
其中有两个疑问语气助词“乎”，
则可以考虑是否该在它们之后各加问号，
表示一个疑问句结束之后的停顿；
其中有三个表示陈述语气的语气词“也”，
则应结合上下文意的分析，
究竟是用于句末表示判断、
确认或者表示申述原因、
进行推理呢，
还是用于句中表示句中停顿呢？
结合辨析它们前面的两个疑问结构
“何故”
和两个连词
“则”
字，
就可以判定前两个
“也”
字是表示申述原因的，
最后一个
“也”
字则明显地表示判断，
它们后面都应该用个句号。
我们再根据连词
“而”
和
“以”
连接前后两个相关的词语
（“争一言以相杀”的“以”不是构成介词结构，而是连接两个相承的动作，所以是连词），
又可以点开两处。
至于
“于义”、
“于其身”
是两个明显的修饰谓语的介词结构，
都可以帮助我们断句和理解文义。
特别是
“今谓人曰”
与
“又曰”
之后的文字前后并列对举，
句式结构基本相似，
更可以帮助我们理解文意和断句。
分析至此，
全段断句和标点也就明确起来了，
那就是：

子墨子曰：
“万事莫贵于义。
今谓人曰：‘予子冠履而断予之手足，子为之乎？’
必不为。
何故？
则冠履不若手足之贵也。
又曰：
‘予子天下而杀予之身，子为之乎？’
必不为。
何故？
则天下不若身之贵也。
争一言以相杀，
是贵义于其身也。”
\newpage